\subsection{从固定阶失效到 Sudakov 重求和}
\label{sec:qcd-sudakov-thrust-applications}

软、共线极限已说明，降低喷注分辨阈值会产生大对数。仅仅指出固定阶展开失效还不能给出可用预测；下一步必须把各阶重复出现的辐射贡献重新组织。一个适合完整演算的观测量是推力：它测量强子末态的动量沿某条最优轴集中到什么程度。

对质心系中所有测得的末态三动量 $\vb{p}_i$，定义
\begin{equation}
 T_{\rm thr}=\max_{|\vb{n}|=1}
 \frac{\sum_i|\vb{p}_i\cdot\vb{n}|}{\sum_i|\vb{p}_i|},
 \qquad \tau_T=1-T_{\rm thr}.
 \label{eq:qcd-thrust-definition}
\end{equation}
最大化得到的 $\vb{n}$ 称为推力轴。无质量末态有 $1/2\leq T_{\rm thr}\leq1$；两个背对背的窄喷注对应 $\tau_T\to0$。增加一个无限软粒子不改变该比值，共线分裂也不改变投影总和，因此推力是红外--共线安全的。有限质量强子的实验定义仍使用实际三动量，不能在精密比较中不加说明地把能量与动量模长互换。

累积分布的精确定义为
\begin{equation}
 \Sigma_T(\tau)=\frac{1}{\sigma_{\rm had}}
 \sum_X\int\dd\sigma_X\,
 \Theta[\tau-\tau_T(X)].
 \label{eq:qcd-thrust-cumulant}
\end{equation}
$\sigma_{\rm had}$ 是同一初态下的包容强子截面，$\dd\sigma_X$ 已含振幅平方、通量与相空间；$\Theta$ 是阶跃函数。这个定义对多粒子末态直接测量 $\tau_T(X)$，尚未假设只辐射一个胶子。精确分布非负、随 $\tau$ 单调增加，并在覆盖全部物理末态时趋于一。

在 $\tau\ll1$ 的两喷注区域，领先幂动力学分别含硬能标 $E_{\rm cm}$、喷注能标 $E_{\rm cm}\sqrt\tau$ 和软能标 $E_{\rm cm}\tau$。在这些尺度都处于微扰区时，可以用硬函数、喷注函数和软函数作系统因子化；它们的重整化群演化组织完整的对数级数\cite{BecherSchwartz2008Thrust}。下面先保留固定 $\alpha_s$ 的领先双对数，逐项算出这一结构最直接的辐射解释。该精度保留 $\alpha_s^n\ln^{2n}\tau$，不包含完整的运行耦合领先对数或次领先对数。

对于相对某条硬夸克方向的小角软胶子，定义能量分数 $z_g=2E_g/E_{\rm cm}$、角变量 $u_\theta=\theta^2/4$。单次发射对推力亏损的领先贡献为 $v=z_gu_\theta$，单条夸克腿的实辐射测度为
\begin{equation}
 \dd\mathcal P_q=\frac{\alpha_sC_F}{\pi}
 \frac{\dd z_g}{z_g}\frac{\dd u_\theta}{u_\theta}.
 \label{eq:qcd-thrust-emission-measure}
\end{equation}
有限角度与非软分裂项会产生较低次对数。只求双对数时，可以把两个变量的上限设为一；改变这些数量级为一的上限会改变单对数与常数项，不能用于确定更高精度的系数。

两条背对背的夸克腿都能辐射。若要求没有单次发射使 $v>\tau$，排除区域中的总辐射权重为
\begin{equation}
 R_T(\tau)=\frac{2\alpha_sC_F}{\pi}
 \int_\tau^1\frac{\dd z_g}{z_g}
 \int_{\tau/z_g}^1\frac{\dd u_\theta}{u_\theta}
 =\frac{\alpha_sC_F}{\pi}\mathcal L_T^2,
 \qquad \mathcal L_T=\ln\frac1\tau.
 \label{eq:qcd-thrust-radiator}
\end{equation}
虚修正和任意多个允许的未分辨实发射共同产生 Sudakov 因子，固定耦合双对数结果为
\begin{equation}
 \Sigma_T^{\rm DL}(\tau)=\ee^{-R_T(\tau)},\qquad
 \frac{\dd\Sigma_T^{\rm DL}}{\dd\tau}
 =\frac{2\alpha_sC_F}{\pi}\frac{\mathcal L_T}{\tau}
 \ee^{-(\alpha_sC_F/\pi)\mathcal L_T^2}.
 \label{eq:qcd-thrust-sudakov}
\end{equation}
在此精度内可以将包容归一化替换成 Born 归一化，因为两者差异从不含双对数的相对 $\order(\alpha_s)$ 开始。重求和并没有消除胶子发射，而是计算观测窗口中始终没有出现可分辨发射的概率。

固定阶展开只保留 $\Sigma_T=1-R_T+\cdots$；当 $R_T>1$ 时会给负的累计率。指数把整组 $R_T^n/n!$ 保留下来，使固定耦合近似的累计率保持在零与一之间。然而概率非负只是一项必要检查，不能据此认定该近似已达到实验精度。在 $\tau\to0$ 时，真实运行耦合会取样低能区域；当 $E_{\rm cm}\tau$ 接近强子化能标，微扰 Sudakov 公式也需要非微扰输入。

\begin{exbox}[大对数怎样使一阶近似失效]
取固定教学耦合 $\alpha_s=0.120$、$C_F=4/3$。分别在 $\tau=0.10$ 和 $0.01$ 计算一阶累计率与双对数指数结果，估计固定阶失效的分辨率。
\textbf{解。} 在 $\tau=0.10$ 时，$R_T\simeq0.270$，一阶累计率为 $0.730$，双对数指数结果为 $0.763$；在 $\tau=0.01$ 时，$R_T\simeq1.080$，一阶结果变为 $-0.080$，指数结果仍为 $0.340$。

固定阶失效的边界由 $R_T\sim1$ 给出，
\begin{equation*}
 \tau\sim\exp\!\left[-\sqrt{\frac{\pi}{\alpha_sC_F}}\right]\simeq0.0119.
\end{equation*}
这解释了小耦合与固定阶可靠性之间为何还隔着一个分辨率条件。若 $E_{\rm cm}=\SI{100}{\giga\electronvolt}$，$\tau=0.01$ 对应的软能标仅为 $\SI{1}{\giga\electronvolt}$；即使硬尺度上的 $\alpha_s$ 很小，也不能据此忽略软区运行与强子化。
\end{exbox}

\begin{figure}[htbp]
 \centering
 \includegraphics[width=0.73\textwidth]{qcd_thrust_sudakov.pdf}
 \caption{固定耦合双对数近似下的推力累计率。虚线是一阶展开，实线是同一双对数级数的指数求和；负的一阶结果标志其失效。横轴只展示两喷注区域，曲线未包含运行耦合、次领先对数与强子化。}
 \label{fig:qcd-thrust-sudakov}
\end{figure}
\figref{fig:qcd-thrust-sudakov} 比较的是同一对数精度的截断与求和。它适合判断截断何时失去层级，定量拟合推力分布还需要更高阶因子化与匹配。

